\subsubsection{Évaluation de séquences}

\paragraph{ }
\subparagraph{Comparaison des fréquences empiriques}
Avant d'évaluer les séquences avec notre modèle de Markov, 
nous avons calculé la fréquence d'apparition simple (modèle d'ordre 0) de chaque 
nucléotide pour les trois séquences fournies :

\begin{center}
\renewcommand{\arraystretch}{1.2}
\begin{tabular}{|l|c|c|c|c|}
    \hline
    \textbf{Séquence} & \textbf{Fréq. A} & \textbf{Fréq. C} & \textbf{Fréq. G} & \textbf{Fréq. T} \\
    \hline
    Séquence 1 & 0.2885 & 0.3495 & 0.2090 & 0.1530 \\ 
    Séquence 2 & 0.3060 & 0.3415 & 0.1970 & 0.1555 \\ 
    Séquence 3 & 0.2845 & 0.3605 & 0.1985 & 0.1565 \\ 
    \hline
\end{tabular}
\end{center}

L'analyse de ce tableau révèle que:  
\textit{Seq1, Seq2 et Seq3 ont des distributions très similaires.}

\subparagraph{Évaluation par le modèle de Markov}
Afin d'évaluer plus finement dans quelle mesure ces séquences sont compatibles avec 
la dynamique de \texttt{seq1}, nous calculons la probabilité d'observer chacune d'entre 
elles sous l'hypothèse qu'elles ont été générées par notre chaîne de Markov.

La vraisemblance d'une séquence $S = (x_1, x_2, \dots, x_L)$ s'écrit :
\begin{equation}
    \Prob(S \mid P) = \pi_\infty(x_1) \prod_{t=1}^{L-1} P_{x_t, x_{t+1}}
\end{equation}

Pour éviter l'underflow numérique dû au produit de milliers de probabilités, 
nous utilisons la log-vraisemblance moyenne par nucléotide :
\begin{equation}
    \frac{1}{L} \log \Prob(S \mid P) = \frac{1}{L} \left( \log(\pi_\infty(x_1)) + \sum_{t=1}^{L-1} \log(P_{x_t, x_{t+1}}) \right)
\end{equation}

Les résultats obtenus pour les trois séquences sont les suivants :
\begin{itemize}
    \item \textbf{Séquence 1} : Log-vraisemblance moyenne = -1.2140
    \item \textbf{Séquence 2} : Log-vraisemblance moyenne = -1.4864
    \item \textbf{Séquence 3} : Log-vraisemblance moyenne = -1.2005
\end{itemize}

\subparagraph{Conclusion : Comparaison et puissance du modèle de Markov}

L'analyse conjointe des fréquences empiriques et des log-vraisemblances 
révèle l'intérêt fondamental de la modélisation par chaînes de Markov. 

Dans un premier temps, l'observation des fréquences d'apparition simples 
montre que les trois séquences possèdent une composition globale quasi-identique 
(environ 29\% de A, 35\% de C, 20\% de G et 16\% de T). Sur la seule 
base de ces comptages, on pourrait conclure que ces trois séquences proviennent 
du même processus stochastique.

Cependant, l'évaluation par la log-vraisemblance met en lumière une réalité différente, 
car elle prend en compte la dynamique temporelle (les probabilités de transition entre 
nucléotides adjacents) :
\begin{itemize}
    \item \textbf{La Séquence 3} obtient une log-vraisemblance moyenne ($-1.2005$) extrêmement proche de celle de la Séquence 1 ($-1.2140$), sur laquelle le modèle a été entraîné. Cela prouve que la Séquence 3 partage la même "signature" de transitions stochastiques que la Séquence 1. Elles ont très certainement été générées par le même processus.
    \item \textbf{La Séquence 2}, en revanche, obtient une log-vraisemblance nettement inférieure ($-1.4864$). L'échelle étant logarithmique sur 2000 caractères, cette différence indique que la probabilité que la Séquence 2 ait été générée par notre matrice $P$ est infime. 
\end{itemize}

Nous pouvons donc conclure que bien que la Séquence 2 utilise les mêmes proportions de 
nucléotides que la Séquence 1, leur ordonnancement est fondamentalement différent. 
Le modèle de Markov réussit là où le simple comptage échoue : il détecte que la 
Séquence 2 est une anomalie générée par un processus 
(ou un modèle de fond biologique) distinct.